% ---
% id: meat-core/frank-wolfe/def:relaxed-feasible-set
% kind: def
% title: "The relaxed feasible set"
% summary: >-
%   At a window ε ≥ 0 the relaxed feasible set keeps positivity and the block
%   structure exact and relaxes the marginal or coupling condition, the declared
%   rows and, with them, the normalisation. At ε = 0 it is the feasible set of
%   the program, in the state form and in the Choi form.
% prominence: standard
% status: implemented
% uses: [meat-core/def:feasible, meat-core/thm:choi]
% mentions: [meat-core/def:untrusted-choi, prop:relaxed-set-properties, foundations/lem:floor-stricter, foundations/req:norm-conversion, meat-core/def:declaration]
% assumes: []
% refs: []
% provenance: newly-derived
% implements: certiqs_key.solvers.frank_wolfe.joint_driver:_check_point
% lean:
% ---

The conditions defining $\Feas$ in \eqref{eq:feas}--\eqref{eq:feas-modes}
are of two kinds. Some compare the variable against \tref{declared data}:
the marginal condition against $\tau_A$, and each relation $\bowtie_k$
against $\Theta_k$ or $c_k$. The others, positive semidefiniteness, the
sector-block structure and the normalisation $\tr\omega=1$, name no
declared datum. A run establishes membership of its iterates only up to
windows on computed residuals, so the set is read at a tolerance; positive
semidefiniteness and the block structure stay exact. The normalisation is
relaxed with the condition that couples the sector blocks because it is an
arithmetic consequence of that condition, and in one mode it is that
condition: in the state form in \tref{trusted-source mode}, $\tr\omega=1$
follows from $\omega_A=\tau_A$; in \tref{untrusted-source mode},
$\sum_s\tr\omega_s=1$ is the coupling condition itself
(Definition~\ref{def:untrusted-choi}\,(b)); in the Choi form,
$\tr[\chi_{\Qsys}(J)]=1$ follows from
$\sum_s\tr_{\Qsys_s}J_s=\idon{A_0}$ (proof of Theorem~\ref{thm:choi}).
Holding the normalisation exact beside a relaxed coupling condition would
exclude every point whose coupling residual $R$ has
$\tr[\iota^{\dagger}\tau_A\iota\,R]\ne0$
(Proposition~\ref{prop:relaxed-set-properties}\,(c$'$)): the generic point
with a nonzero residual, the residual being the drift over the convex
combinations that built the iterate, and a nonzero residual is what the
coupling check of a run accepts in the normal case. The dimension factors
$d_A$ and $d_{A_0}$ of the normalisation clauses of
Definition~\ref{def:relaxed-feasible-set} are written into them as \xref{foundations/req:norm-conversion} requires.

\begin{definition}\label{def:relaxed-feasible-set}
For $\varepsilon\ge0$ the
\term[The feasible set read at a window: positivity and the sector-block structure kept exact, the marginal or coupling condition, the declared constraint rows and the normalisation each relaxed at the window.]{relaxed feasible set}
$\Feas(\varepsilon)$ is the set of sector-block-diagonal
$\omega\in\Pos{A\Qsys}$ satisfying, in trusted-source mode,
\begin{equation}\label{eq:feas-eps}
  \big\|\omega_A-\tau_A\big\|_{\max}\le\varepsilon,
  \quad
  \big|\tr\omega-1\big|\le d_A\varepsilon,
  \quad
  \begin{cases}
  \lmin\big(\Theta_k-A_k(\omega)\big)\ge-\varepsilon, & \bowtie_k\ =\ \preceq,\\
  \lmin\big(A_k(\omega)-\Theta_k\big)\ge-\varepsilon, & \bowtie_k\ =\ \succeq,\\
  \big\|A_k(\omega)-\Theta_k\big\|_{\max}\le\varepsilon, & \bowtie_k\ =\ =,
  \end{cases}
\end{equation}
where $A_k(\omega)\coloneqq\tr_{\Qsys}[(\idon{A}\otimes\hat O_k)\,\omega]$
is the left-hand side of the $k$-th declared relation of
\eqref{eq:feas-modes}, and, in untrusted-source mode, the normalisation
$\big|\sum_s\tr\omega_s-1\big|\le\varepsilon$ together with the scalar
forms
$a_k(\omega)\le c_k+\varepsilon$, $a_k(\omega)\ge c_k-\varepsilon$,
$|a_k(\omega)-c_k|\le\varepsilon$, with
$a_k(\omega)\coloneqq\tr[(\idon{A}\otimes\hat O_k)\,\omega]$ and no
marginal condition. In the Choi form $\Feas^{J}(\varepsilon)$ is the set of
$(J_s)_s$ with $J_s\in\Pos{A_0\Qsys_s}$,
$\big|\tr[\chi_{\Qsys}(J)]-1\big|\le d_{A_0}\varepsilon$,
$\big\|\sum_s\tr_{\Qsys_s}J_s-\idon{A_0}\big\|_{\max}\le\varepsilon$ and
the same relaxation of the $\bowtie_k$ rows.

The normalisation clause of each form is implied by the clause written
before it, by the marginal in trusted-source mode and by the coupling
residual in the Choi form, and is written out only so that the set can be
read without re-deriving it; Proposition~\ref{prop:relaxed-set-properties}\,(c$'$)
proves both implications and derives the two dimension factors. In
untrusted-source mode there is no clause before it: there the
normalisation \emph{is} the relaxed coupling condition
(Definition~\ref{def:untrusted-choi}\,(b)), which is why it carries no
dimension factor.

Which predicate is which is part of the definition and is never
converted: an inequality between operators is relaxed by a \emph{spectral
floor} $\lmin(\cdot)\ge-\varepsilon$, an operator equality by an
\emph{entrywise} ball $\|\cdot\|_{\max}\le\varepsilon$, a scalar equality
by an interval, and these are different quantities against the same number
(\xref{foundations/lem:floor-stricter}: the entrywise reading and the
spectral floor). The factors $d_A$ and $d_{A_0}$ are not such a
conversion: each is the dimension factor of an explicit trace bound
derived in Proposition~\ref{prop:relaxed-set-properties}\,(c$'$), written
into the clause, and neither reads an entrywise defect as a spectral one.

At $\varepsilon=0$ both sets reduce to \eqref{eq:feas} and
\eqref{eq:choi-feas} respectively: each normalisation clause is then the
exact $\tr\omega=1$, respectively $\tr[\chi_{\Qsys}(J)]=1$, which the
coupling condition already forces, so re-imposing it changes nothing. We
write $\Feas$ for $\Feas(0)$ and $\Feas^{J}$ for $\Feas^{J}(0)$
throughout.

\emph{In the general format} of Definition~\ref{def:declaration} both sets
are read on lists over the \tref[joint sector]{joint sectors}, with every
line of \eqref{eq:feas-general}, respectively of the general-format Choi
set of Theorem~\ref{thm:choi}, relaxed by the predicate its relation calls
for, as in \eqref{eq:feas-eps}, the \tref{party-local} rows included. The
marginal condition and the coupling condition are then read per Alice
sector:
$\|\sum_s\tr_{\Qsys_s}\omega^{(\mathsf t,s)}-\tau^{(\mathsf t)}_A\|_{\max}
\le\varepsilon$, respectively
$\|R_{\mathsf t}\|_{\max}\le\varepsilon$ with
$R_{\mathsf t}\coloneqq\sum_s\tr_{\Qsys_s}J_{(\mathsf t,s)}
-\idon{A_{0,\mathsf t}}$, for every $\mathsf t\in\mathbb{T}$. The residual
of the single marginal clause of \eqref{eq:feas-eps}, respectively of the
single coupling clause of the Choi form, is then the block-diagonal operator
$\omega_A-\tau_A=\bigoplus_{\mathsf t}(\sum_s\tr_{\Qsys_s}
\omega^{(\mathsf t,s)}-\tau^{(\mathsf t)}_A)$, respectively
$\bigoplus_{\mathsf t}R_{\mathsf t}$ with
$J_s\coloneqq\bigoplus_{\mathsf t}J_{(\mathsf t,s)}$, whose entrywise
maximum is the largest over $\mathsf t$ of the blocks' ones; so the clauses
per Alice sector together are that single clause. The normalisation clauses
are those of the first paragraph of this definition. At $\varepsilon=0$ the
two sets are $\Feas$ of \eqref{eq:feas-general} and the general-format
$\Feas^{J}$ of Theorem~\ref{thm:choi}.
\end{definition}
